\documentclass[10pt]{article}

\usepackage[margin=0.65in]{geometry}
\usepackage{amsmath,amssymb,amsthm}
\usepackage{booktabs}
\usepackage{hyperref}
\usepackage{microtype}

\setlength{\parskip}{0.1em}
\setlength{\parindent}{1.1em}
\linespread{0.92}

\newtheorem{proposition}{Proposition}
\newtheorem{assumption}{Assumption}

\title{\vspace{-1.6em}Real Determinacy Tests for DSGE Estimation}
\author{Matyas Farkas}
\date{\vspace{-0.9em}\today}

\begin{document}
\maketitle
\vspace{-1.1em}

\begin{abstract}
\noindent
DSGE estimation usually rejects parameter draws for which the linearized model
has no unique stable solution. This paper shows that one common rejection rule
is too strict: a draw should not be rejected only because a complex-QZ
calculation returns a complex-valued policy matrix. In a real linear
rational-expectations model, stable complex roots come in conjugate pairs and
describe damped cycles in real variables. If the stable-root count and the
state-to-jump mapping condition hold, the policy rule is real. In one million
loose-prior draws from the benchmark HLT model, the complex-QZ rule rejects
464,685 draws that real QZ accepts with median residual
$8.8\times10^{-14}$. A Cayley real-part rule matches real QZ in this audit;
doubling accepts 84,606 draws that real QZ classifies as non-unique.
\end{abstract}
\vspace{-0.4em}
\noindent\textbf{Keywords:} DSGE estimation; determinacy; generalized Schur;
Bayesian posterior support; GPU computing.

\section{Introduction}

A dynamic stochastic general equilibrium (DSGE) model describes how economic
variables move over time in response to shocks. In Bayesian DSGE estimation,
each parameter draw must pass a solution test before its likelihood is used.
The standard first-order test is determinacy: the linearized model must have
one stable equilibrium path for each initial state. If the test is wrong, the
posterior is wrong because valid parameter draws are either excluded or invalid
draws are included.

Formally, the posterior density is proportional to
\[
    p(\theta \mid Y) \propto L(Y \mid \theta)\pi(\theta)
    \mathbf{1}\{\theta \in \Theta_D\},
\]
where $\theta$ is the parameter vector, $Y$ is the data, $L(Y\mid\theta)$ is
the likelihood, $\pi(\theta)$ is the prior, and $\Theta_D$ is the set of
determinate parameter values. The indicator is not a harmless numerical detail.
It defines the parameter region over which the posterior is computed.

This paper proves a narrow point about that indicator. A complex-valued
intermediate result from a complex-QZ calculation is not enough to reject a
draw. QZ, also called the generalized Schur decomposition, is the standard
algorithm for separating stable and unstable generalized eigenvalues in linear
rational-expectations models (Blanchard and Kahn, 1980; Klein, 2000; Sims,
2002). Complex arithmetic can appear inside that calculation even when the
economic solution is real. The valid rejection criteria are the stable-root
count, the ability to write jump variables as a function of states, and the
residual of the recovered policy rule.

The distinction matters for GPU code. Full generalized QZ is difficult to run
efficiently in JAX on GPUs. Faster alternatives, such as Cayley transforms or
doubling iterations, are useful only if they make the same accept/reject
decision as real QZ on the parameter region used in estimation. The audit below
shows that the old complex-QZ rule falsely rejects many valid draws in the HLT
benchmark model. It also shows that doubling is too loose as a final
determinacy test.

\section{Objects and Definitions}

After linearization, a DSGE model can be written as a generalized eigenvalue
problem. Let $(D(\theta),E(\theta))$ be the real matrix pencil at parameter
value $\theta$. Its generalized eigenvalues solve
\[
    \det(D(\theta)-\lambda E(\theta))=0.
\]
Finite roots with $|\lambda|<1$ are stable: repeated application of the
linearized dynamics makes them decay. Infinite generalized roots can also
appear when the denominator in the generalized eigenvalue ratio is zero. They
are not counted as finite stable roots.

The variables are ordered into states and jumps. State variables are fixed when
current choices are made, for example capital or lagged shocks. Jump variables
can adjust immediately, for example prices or interest rates. Let $n_s$ be the
number of state variables. A first-order solution is a matrix $G$ such that
\[
    y_t = Gx_t,
\]
where $x_t$ is the vector of states and $y_t$ is the vector of jump variables.

The stable generalized eigenvectors span the stable deflating subspace. Write
one basis for that subspace as
\[
    Z =
    \begin{bmatrix}
        Z_x \\
        Z_y
    \end{bmatrix},
\]
where $Z_x$ corresponds to states and $Z_y$ corresponds to jumps. The graph
condition says that $Z_x$ is nonsingular. In plain language, the stable
subspace must give exactly one jump-vector choice for each state vector. When
the graph condition holds, the policy matrix is $G=Z_yZ_x^{-1}$.

The determinacy set used in estimation is therefore
\[
    \Theta_D =
    \left\{
        \theta :
        \#\{\text{finite }\lambda_i(\theta):|\lambda_i(\theta)|<1\}=n_s,\;
        Z_x(\theta)\text{ is nonsingular},\;
        r(\theta)\leq \tau
    \right\}.
\]
Here $r(\theta)$ is the relative residual of the recovered policy rule, and
$\tau$ is the numerical tolerance. A rule for accepting or rejecting draws is
valid on a parameter region only if it gives the same indicator as
$\mathbf{1}\{\theta\in\Theta_D\}$ on that region.

\begin{assumption}[real regular pencil]
For the parameter vector under consideration, $D$ and $E$ are real, the pencil
$(D,E)$ is regular, and no finite generalized eigenvalue lies on the unit
circle.
\end{assumption}

\begin{assumption}[state-to-jump graph]
The stable deflating subspace has dimension $n_s$. For any full-rank basis
$Z=(Z_x',Z_y')'$ of this subspace, the state block $Z_x$ is nonsingular.
\end{assumption}

\section{Claim and Proof}

\begin{proposition}[complex stable roots imply real damped cycles]
Under Assumptions 1 and 2, the first-order policy rule
\[
    y_t = Gx_t,\qquad G=Z_yZ_x^{-1},
\]
is real. Complex finite generalized eigenvalues with modulus below one do not
imply instability. If a complex-QZ implementation returns a policy matrix with
non-negligible imaginary entries despite the correct stable-root count and
state-to-jump graph, the complex part is a numerical extraction problem, not by
itself evidence that $\theta\notin\Theta_D$.
\end{proposition}

The key idea is that a real model maps real states into real choices. A complex
root is a convenient way to describe a cycle. Its conjugate root appears with
it. Together the two roots describe a real two-dimensional rotation that shrinks
over time when their modulus is below one.

\begin{proof}
Because $D$ and $E$ are real, conjugation preserves generalized eigenpairs. If
$Dv=\lambda Ev$, then $D\bar v=\bar\lambda E\bar v$. Thus non-real generalized
eigenvalues and their eigenvectors occur in conjugate pairs. The stable set
$\{\lambda:|\lambda|<1\}$ is closed under conjugation, so the stable deflating
subspace is also closed under conjugation.

Any conjugation-closed complex subspace has a real basis. For each non-real
vector $v=a+ib$, replace the pair $(v,\bar v)$ with the real vectors
$a=\operatorname{Re}v$ and $b=\operatorname{Im}v$. These two real vectors span
the same plane as the complex pair. Repeating this replacement gives a real
full-rank basis $Z=(Z_x',Z_y')'$ for the stable subspace.

Assumption 2 says that $Z_x$ is nonsingular. The stable subspace is therefore
the graph of
\[
    G=Z_yZ_x^{-1}.
\]
Both matrices in this expression are real, so $G$ is real. The result does not
depend on the particular basis. If $\tilde Z=ZT$ is another full-rank complex
basis for the same subspace, then
\[
    \tilde Z_y\tilde Z_x^{-1}
    = Z_yT(Z_xT)^{-1}
    = Z_yZ_x^{-1}.
\]
In exact arithmetic, every correct basis gives the same real policy matrix.

Finally, a complex stable root $\lambda=a+ib$ corresponds to a real $2\times2$
block with eigenvalues $a\pm ib$. Powers of that block decay at rate
$|\lambda|$. The path can oscillate, but it converges. This is stability in the
real state space.
\end{proof}

\section{Audit on a Benchmark DSGE Model}

The proof says that complex stable roots are not a valid rejection reason. The
next question is numerical: how often does the old rejection rule matter? I
answer this with one million loose-prior parameter draws from the HLT benchmark
model. HLT is used here as the model label in the replication code; the proof
above does not rely on special features of this model.

Each draw is classified four ways. Real QZ is the reference calculation.
Complex QZ is the old rule that rejects complex-valued extracted policy
matrices. The Cayley real-part rule is a faster screen designed for batched
linear algebra. Doubling is a fast iterative solver. The audit file is
\texttt{benchmarks/results/qz\_real\_cayley\_hlt\_loose\_1m\_stream.json}.

\begin{table}[h!]
\centering
\caption{Accept/reject decisions in one million loose-prior HLT draws}
\label{tab:audit}
\begin{tabular}{lrr}
\toprule
Classification event & Count & Share \\
\midrule
Real QZ accepts: unique stable solution & 702,789 & 70.3\% \\
Real QZ rejects: indeterminate & 256,850 & 25.7\% \\
Real QZ rejects: no stable solution & 40,361 & 4.0\% \\
Complex-QZ accepts: unique stable solution & 238,104 & 23.8\% \\
Complex-QZ rejects as ``complex solution'' & 464,685 & 46.5\% \\
Real QZ accepts those ``complex'' cases & 464,685 & 46.5\% \\
Cayley real-part disagrees with real QZ & 0 & 0.0\% \\
Doubling accepts a real-QZ non-unique case & 84,606 & 8.5\% \\
\bottomrule
\end{tabular}
\vspace{0.25em}

\footnotesize
Notes: The median real-QZ residual among accepted draws is
$8.8\times10^{-14}$; the maximum is $1.0\times10^{-11}$. The Cayley
real-part rule matches real QZ in this audit. This is evidence on the tested
HLT support, not a theorem for every DSGE pencil.
\end{table}

Table \ref{tab:audit} shows that the old complex-QZ rule changes the accepted
parameter region. It rejects 464,685 draws because the extracted policy matrix
is complex. Real QZ accepts all of these draws and rejects none of the draws
that complex QZ labels unique. The old rule therefore removes many real,
stable, unique solutions.

The table also shows why fast solvers still need verification. Doubling finds
all real-QZ unique cases, but it also accepts 84,606 cases that real QZ
classifies as non-unique. Doubling can be useful as a fast intermediate step,
but it should not be the final determinacy test. In this audit, the Cayley
real-part rule agrees with real QZ on the accept/reject decision.

\section{Implications for Estimation}

The practical rule is simple. First, count the finite stable roots. Second,
check that the stable subspace defines one jump vector for each state vector.
Third, recover a real policy matrix and verify the residual. A complex
intermediate matrix should trigger a real reconstruction or a fallback to real
QZ. It should not by itself reject the parameter draw.

This rule is useful for GPU estimation. JAX can run Cayley-style screens and
batched residual checks on the GPU. Real QZ should remain the reference
validator until a reliable GPU implementation of generalized QZ is available.
Doubling can still speed up the pipeline, but only behind a verification step.

\section{Conclusion}

Complex roots in a real DSGE model describe damped cycles, not complex-valued
economic choices. The valid determinacy test is the stable-root count, the
state-to-jump mapping condition, and the residual of the recovered policy
rule. In the HLT audit, replacing the old complex-QZ rejection rule with real
QZ changes the number of accepted draws from 238,104 to 702,789. This change is
large enough to alter posterior estimates.

\section*{Acknowledgments and AI Use Disclosure}

The author credits OpenAI's language models for drafting assistance, editing
suggestions, and help converting computational audit results into manuscript
text. The author remains responsible for all claims, proofs, numerical
interpretations, and citations.

\begin{thebibliography}{9}
\setlength{\itemsep}{0pt}
\small

\bibitem{bk1980}
Olivier Jean Blanchard and Charles M. Kahn. 1980.
The solution of linear difference models under rational expectations.
\emph{Econometrica} 48(5): 1305--1311.

\bibitem{klein2000}
Paul Klein. 2000.
Using the generalized Schur form to solve a multivariate linear rational
expectations model.
\emph{Journal of Economic Dynamics and Control} 24(10): 1405--1423.

\bibitem{sims2002}
Christopher A. Sims. 2002.
Solving linear rational expectations models.
\emph{Computational Economics} 20: 1--20.

\bibitem{lapack1999}
LAPACK Working Group. 1999.
\emph{LAPACK Users' Guide}. Third edition. SIAM.

\end{thebibliography}

\end{document}
